\honest{When Anthropomorphism Became Stupid}{Eliezer Yudkowsky}{2008}

Most things in the universe do not have minds. Earlier cultures thought that trees, rocks, streams and hills had spirits, and I say this was not stupid. Human flesh thinks, so why not wood? My muscles move at my will; when the river floods the tribe's gathering-place, why not think the river is angry, as you would of a fist in your face? \nb{This is Hume's account of popular religion, in \textsc{The Natural History of Religion} (1757): ``There is an universal tendency among mankind to conceive all beings like themselves.'' Hume also said what would cure it, knowledge of ``the minute parts'' of bodies, which is close to the cure this post reaches.}

The belief only looks stupid because it is not our tribe's. To make it obvious that wood does not think, I say, ``you have to belong to a culture with microscopes.'' \nb{Aristotle, with no microscope, held that plants have no power of thinking, reasoning from what plants are observed to do and from a theory of perception.}

Aristotle thought the brain cooled the blood; Egyptian embalmers threw it away; Alcmaeon's view that it was the seat of intelligence was only a guess. When did it stop being a guess? ``I do not know enough history to answer this question.'' Maybe when someone found that cutting the nerves to the brain blocked movement and sensation. Even then, what moved through the nerves could be a mysterious spirit that wood and water share. \nb{Galen cut nerves and the spinal cord of animals in the second century, and held that a ``psychic pneuma'' filled the brain and nerves. The post's guess and its answer to the guess both fit him.} I searched for the moment someone took the brain's tangle of neurons for information processing, and ``I haven't had much luck.''

I put the beginning of the change at that tangle, at Darwin, and at the idea of cognition as computation. After them you can say why a tree, a river or an atom has no mind, and why a puppy is not doing moral reasoning, according to ``Our current theories of evolutionary psychology,'' which I do not cite. And you can ask a rock where it would get ``the intentions to want to roll downhill, as Aristotle once thought.'' \nb{Aristotle wrote that falling bodies do not move themselves, since self-motion ``is a characteristic of life and peculiar to living things.''}

I end with Zhuangzi and Huizi arguing about whether a man can know what fish enjoy, and add: ``Now we know.''
